%% kapa: exam=1 round=31 year=2020 period=1 form=A subject=real_estate subject_ko=부동산학원론 seq=24 number=104
%% source: Q-Net 감정평가사 기출문제 (cst00309) artlSeq=5207219 file=extracted/1008/5207219/A형/A형/2020년도 제31회 감정평가사 1차 시험_1교시_A형.pdf
%% conversion: kordoc 4.24.51 -> tools/gen1.py (auto); review=proofread-sonnet; flags=none
\begin{question}{104}
저당대출의 상환방식에 관한 설명으로 옳은 것은?
\begin{choices}
  \item 원금균등분할상환(CAM) 방식의 경우, 원리금의 합계가 매기 동일하다.
  \item 원리금균등분할상환(CPM) 방식의 경우, 초기에는 원리금에서 이자가 차지하는 비중이 높으나, 원금을 상환해 가면서 원리금에서 이자가 차지하는 비중이 줄어든다.
  \item 다른 조건이 일정하다면, 대출채권의 듀레이션(평균 회수기간)은 원리금균등분할상환(CPM) 방식이 원금균등분할상환(CAM) 방식보다 짧다.
  \item 체증분할상환(GPM) 방식은 장래 소득이 줄어들 것으로 예상되는 차입자에게 적합한 대출방식이다.
  \item 거치식(Interest-only Mortgage) 방식은 대출자 입장에서 금리수입이 줄어드는 상환방식으로, 상업용 부동산 저당대출보다 주택 저당대출에서 주로 활용된다.
\end{choices}
\end{question}
\answer{2}
